
\subsubsection{6.1 Interpretation of Findings}

\textbf{Adversarial construction method determines calibration outcome.} Human-adversarial NLI examples (AdvGLUE MNLI: $\Delta ECE$ = +0.071) produce calibration degradation; model-in-loop adversarial NLI examples (ANLI R1/R2: $\Delta ECE$ < 0) produce calibration improvement. This distinction---absent in prior adversarial robustness work, which pools all ''adversarial'' conditions---is the key methodological contribution of this experimental design.

The distinction also provides a bridge between these NLP findings and the vision-domain calibration literature. Minderer et al. [2021] demonstrated that image corruption consistently degrades calibration. Image corruptions are not model-aware: they do not select examples that the model was already failing on. AdvGLUE human-adversarial examples share this property. ANLI model-in-loop examples do not: they were selected to cause errors in the model, so the model is already appropriately uncertain, producing calibration improvement rather than degradation. Model-awareness of the adversarial construction method appears to be the critical variable explaining the NLP versus vision discrepancy.

\textbf{RLHF alignment exhibits a benchmark-type-conditional calibration effect.} The alignment tax on AdvGLUE (chat worse than base: $\Delta \Delta ECE = -0.026)$ and the consistent RLHF moderation on ANLI (100\% of rounds, $\Delta \Delta ECE$ up to +0.147) reflect the same underlying interaction. RLHF-aligned models trained to produce confident, helpful outputs show reduced calibration degradation when adversarial examples probe uncertainty (ANLI model-in-loop) but may show increased degradation when adversarial examples exploit decisive wrong answers (AdvGLUE human-adversarial). This interpretation is consistent with the data but awaits direct experimental confirmation.

The benchmark-type $\times$ RLHF interaction has practical implications: calibration evaluations conducted on clean benchmarks (as in Kadavath et al. [2022]) do not predict adversarial calibration behavior, and benchmark construction method must be specified when reporting RLHF calibration results.

\textbf{Task type is an independent calibration moderator.} AdvGLUE QQP (binary, human-adversarial) shows $\Delta ECE = -0.029$ despite having the same construction method as AdvGLUE MNLI ($\Delta ECE$ = +0.071). Distributing logits across three options (NLI) versus two options (binary) produces meaningfully different susceptibility to adversarial calibration effects.

\subsubsection{6.2 Limitations}

\textbf{L1: Single-model evaluation for primary calibration experiments.} Core calibration results (RQ1--RQ3) use Llama-2-7b-hf only. The planned 4-model grid (Llama-2-7b-hf, Llama-2-7b-chat, Llama-2-13b-chat, Mistral-7B-instruct) was not completed. Consequently, the pre-specified prediction that $\geq 60$\% of model $\times$ task cells would show $\Delta ECE$ > 0.05 is evaluated on only 5 cells from 1 model---statistically underpowered for a universal threshold claim. Results are framed as a pilot study establishing the conditional structure within a single model family. The RLHF moderation finding (RQ4) extends to a model pair within the same family.

\textbf{L2: Shared clean baseline for ANLI.} All ANLI adversarial cells (R1/R2/R3) share one GLUE MNLI clean baseline (n=200, ECE=0.279). This creates cross-cell correlation in ANLI $\Delta ECE$ values. Llama-2-7b-hf's clean MNLI ECE (0.279) is above the Kadavath et al. [2022] range for clean LLMs (0.05--0.15), reflecting weak NLI accuracy (0.365). ANLI R1/R2 $\Delta ECE$ values may be partially driven by an inflated clean baseline. However, AdvGLUE MNLI uses the same baseline and shows the expected positive $\Delta ECE$, so baseline inflation does not fully explain the ANLI calibration improvement---particularly since accuracy on ANLI R1 is actually higher than on clean MNLI (0.380 vs. 0.365), consistent with an adaptive uncertainty explanation.

\textbf{L3: BBH-MC commonsense reasoning not tested.} BIG-Bench Hard multiple-choice adversarial variants were planned but excluded for scope reasons. All claims are bounded to NLI and binary paraphrase classification. NLI is the primary adversarial NLP benchmark domain (both AdvGLUE and ANLI are NLI-centered); the task-type finding (NLI versus binary) is internally complete.

\textbf{L4: RLHF moderation conditionality discovered post-hoc.} The benchmark-type $\times$ RLHF interaction was not predicted in the original H-C1 design; it emerged when H-C1-V2 reported the AdvGLUE reversal as an unexpected finding. Further experimental validation---testing additional RLHF-aligned models, controlling for prompt format---is needed.

\textbf{L5: Pre-specified quantitative thresholds not met.} Two of the pipeline's quantitative predictions were not satisfied at their original thresholds. Mean confidence on incorrect adversarial predictions was 0.616---below the originally hypothesized $\geq 0.70$ threshold (informed by vision-domain analogues). Only 1 of 5 adversarial cells exceeded $\Delta ECE$ > 0.05; the pre-specified criterion was $\geq 60$\% of cells. These results indicate that adversarial calibration degradation in LLMs is smaller in magnitude than vision-domain analogues suggested. The conditional structure remains valid at observed magnitudes.

\subsubsection{6.3 Broader Implications}

Practitioners who test model reliability under adversarial conditions by measuring accuracy robustness alone are missing the calibration dimension: a model can be simultaneously more accurate \textit{and} less calibrated under adversarial conditions, or vice versa. The ANLI R1 result (accuracy improves from 0.365 to 0.380; ECE improves from 0.279 to 0.239) illustrates that accuracy and calibration do not covary in a predictable way under adversarial conditions. Adversarial calibration audit ($\Delta ECE$ measurement) provides a complementary reliability check.

The finding that RLHF alignment can exacerbate calibration degradation on static human-adversarial benchmarks is a potential concern for high-stakes deployment of RLHF-aligned models in NLI or classification settings where adversarial examples were designed to exploit base model reasoning shortcuts.

